\section{Introduction}
\label{sec:intro}

A resilient timing user---a telecom primary reference time clock, a power-grid
phasor unit, a financial timestamping gateway, or a navigation receiver serving
time---rarely fails all at once. It \emph{degrades}: a multi-source, cross-checked
regime, in which the receiver runs a solution-separation integrity monitor over
several independent time sources and reports a timing protection level (TPL) on
its estimate of $|t-\utck{}|$, gives way, on the loss of those sources, to a
single disciplined oscillator running open-loop---\emph{holdover}. Both regimes
have a protection-level literature. What has not been written down is a single PL
valid \emph{across} the transition between them. This paper writes that object
down, proves the one structural property that makes it safe at the seam, and
anchors its holdover leg to a real multi-year metrology series.

\paragraph{The seam is where the safety is lost.}
The danger is specific and, once named, obvious. In the cross-checked regime the
TPL carries the fault-mode overbounds of the integrity monitor---the possibility
that one of the sources is faulty by up to its minimum detectable bias. In pure
holdover there are no cross-check sources and therefore no such fault hypotheses;
a naively re-derived holdover PL at coast time $\tau=0$ can be \emph{smaller} than
the TPL that was being reported an instant earlier. A PL that jumps \emph{down} at
the moment protection is most fragile---the instant the last independent source is
lost---is not merely inelegant; it is unsafe, because the integrity risk it
advertises is momentarily understated. Any composed PL must therefore be
constructed so that it cannot fall across handover. That is the property we prove.

\paragraph{The prior art, and where our contribution begins.}
We bound our novelty up front, because every ingredient here has a parent.
Solution separation and its multiple-hypothesis budget allocation (MHSS) are
aviation ARAIM, cited from Blanch \citep{blanch2015} and from Joerger
\citep{joerger2014}; we specialize their algebra to a degenerate scalar geometry
and claim no new estimator. A receiver-autonomous timing protection level for
GNSS already exists---Gioia's $\TPL=\mathrm{WATP}+\TPL_{\mathrm{noise}}$
\citep{gioia2023}---but it is GNSS-only and single-regime, with no holdover
composition. The holdover leg is built on the standard power-spectral-density to
phase-error-variance clock model of Zucca and Tavella \citep{zucca2005}, from
which the white-FM and random-walk-FM phase-variance growth laws $\tau,\tau^{3}$
follow (the drift $\tau^{5}$ leg being the standard random-run extension); we cite
the model and add nothing to it. The
fault-tolerant selection of a trustworthy time source among redundant domains is
standardized in IEEE~802.1ASed \citep{ieee8021ased}, but the standard
\emph{selects}; it does not emit a protection level or an integrity-risk budget.
The resilient-PNT user-equipment framework IEEE~P1952 \citep{ieeep1952} specifies
a qualitative holdover-drift-\emph{consistency} check at its ``Verify'' level, but
leaves the quantitative PL/IR layer open. Closest of all is the recent conditional
timing protection level of Baweja \citep{baweja2026tpl}, which composes a detector
floor with a coast-over-latency term for a \emph{single} GNSS-disciplined source
and proves that no \emph{unconditional} holdover bound exists. That result is the
$N=1$, single-transition special case of the present construction; we cite it and
generalize it, and we never restate it as new.

\paragraph{What is genuinely ours.}
The contribution is narrow and, we hope, useful.
\begin{enumerate}
  \item \textbf{A scalar timing TPL as a rank-one ARAIM specialization}
  (\Cited{} parent, \Proven{} specialization; Section~\ref{sec:scalar}). When
  every source estimates the \emph{same} scalar $\utck{}$ offset the geometry
  matrix degenerates to $H=\bm 1_N$, so the geometry-dilution structure of
  position ARAIM vanishes and the PL is dominated by the worst
  single-source-\emph{exclusion} subset's noise overbound plus its traceability
  bias: the \emph{bias model}, not geometry, drives the PL. This is Blanch and
  Joerger specialized to a degenerate geometry, not a new theorem.
  \item \textbf{The composed, $\Pzero$-seeded ride-through PL and its
  lower-semicontinuity} (\Proven{}; Section~\ref{sec:composition}). We define
  $\PL(\tau)=\max\big(\TPL_{\mathrm{handover}},\HPL(\tau)\big)$ and prove
  $\PL(0^{+})\ge\PL(0^{-})$: lower-semicontinuous and non-decreasing across
  handover, with no downward jump. Seeding the coast with the handover-state
  covariance $\Pzero$ is proven to be a safety requirement.
  \item \textbf{A law-of-the-iterated-logarithm envelope for the holdover leg}
  (\Proven{}, asymptotic; Section~\ref{sec:lil}). We prove the impossibility of a
  finite deterministic worst-case holdover (the Brownian sup diverges a.s.) and
  the consistency of the IR-allocated $\HPL$ with the almost-sure envelope as
  $\IR\to0$, turning Baweja's impossibility into a positive theorem---while
  stating plainly that the LIL is asymptotic and certifies neither a finite-$t$
  value nor the \num{e-7} quantile.
  \item \textbf{Exactly one externally validated pipeline} (\Validated{};
  Section~\ref{sec:validation}). The holdover leg's running-max envelope
  overbound-covers a real BIPM Circular-T $[\utck{}-\mathrm{UTC(USNO)}]$ series on
  a disjoint segment for every multi-year lag $\tau\ge\SI{90}{\day}$, with the
  short-lag boundary disclosed as \Modelled{}.
\end{enumerate}

\paragraph{Honesty scope.}
We separate four kinds of statement and label them at every occurrence. The MHSS
solution-separation machinery, the clock PSD-to-variance model, the fault-tolerant
selection and resilient-PNT standards, the delay-attack constraint, and the
single-source conditional TPL are \Cited{}. The rank-one specialization theorem,
the lower-semicontinuity of the composed PL, the $\Pzero$-seeding safety property,
and the LIL impossibility-and-consistency statements are \Proven{} algebra, each
machine-checked by the open engine's test suite. Exactly one row is externally
\Validated{}: the holdover-leg multi-year ($\tau\ge\SI{90}{\day}$) overbound
coverage against the real BIPM series. Everything else---the composition at the
system level, the heterogeneous multi-source budget, and every scenario magnitude
and count read as an operational tolerance---is \Proven{} algebra over
\Modelled{} scenarios. The word \Validated{} is never attached to the composition,
the heterogeneous budget, or any end-to-end system claim; those have no external
oracle and are not validated. There is no flight data, no certified alert limit,
no endorsement by any agency, and no technology-readiness claim; every number is
regenerated from the open engine by the commands of
Section~\ref{sec:availability}.
